\documentclass[12pt,reqno]{amsart}

% ============================================
% PACKAGES
% ============================================
\usepackage[utf8]{inputenc}
\usepackage[T1]{fontenc}
\usepackage{amsmath,amssymb,amsthm}
\usepackage{mathtools}
\usepackage{geometry}
\usepackage{hyperref}
\usepackage{booktabs}
\usepackage{graphicx}
\usepackage{xcolor}
\usepackage{cite}

\geometry{letterpaper,margin=1in}

% ============================================
% HYPERLINKS
% ============================================
\hypersetup{
    colorlinks=true,
    linkcolor=blue,
    citecolor=blue,
    urlcolor=blue,
    pdftitle={The Divergence Threshold of the Collatz Map},
    pdfauthor={Luciano Benjamin Nieto}
}

% ============================================
% THEOREM ENVIRONMENTS
% ============================================
\theoremstyle{plain}
\newtheorem{theorem}{Theorem}[section]
\newtheorem{lemma}[theorem]{Lemma}
\newtheorem{proposition}[theorem]{Proposition}
\newtheorem{corollary}[theorem]{Corollary}
\newtheorem{conjecture}[theorem]{Conjecture}
\newtheorem{hypothesis}[theorem]{Hypothesis}

\theoremstyle{definition}
\newtheorem{definition}[theorem]{Definition}
\newtheorem{remark}[theorem]{Remark}
\newtheorem{example}[theorem]{Example}

% ============================================
% CUSTOM COMMANDS
% ============================================
\newcommand{\N}{\mathbb{N}}
\newcommand{\Z}{\mathbb{Z}}
\newcommand{\R}{\mathbb{R}}
\newcommand{\E}{\mathbb{E}}
\newcommand{\PP}{\mathbb{P}}
\newcommand{\nuTwo}{\nu_2}
\newcommand{\logtwo}{\log_2}
\newcommand{\logfour}{\log_4}
\newcommand{\fP}{f_P}
\newcommand{\fPstar}{f_P^*}
\newcommand{\muP}{\mu_P}
\newcommand{\muN}{\mu_N}
\newcommand{\PhiDrift}{\Phi}
\newcommand{\Vmetric}{V_{4/3}}
\DeclareMathOperator{\Orb}{\mathcal{O}}

% ============================================
% METADATA
% ============================================
\title[The Divergence Threshold of the Collatz Map]{%
The Divergence Threshold of the Collatz Map:\\
A Conditional Closed-Form Necessary Condition%
}

\author{Luciano Benjam\'{i}n Nieto}
\address{Independent Researcher, General Alvear, Mendoza, Argentina}
\email{lucianobenjaminnieto@gmail.com}

\date{July 25, 2026}
\subjclass[2020]{Primary 11B83; Secondary 37A05, 11K31, 37P05}
\keywords{Collatz conjecture, $2$-adic dynamics, drift analysis, divergence threshold, ergodic theory}

\begin{document}

\begin{abstract}
We derive an exact closed-form candidate threshold for divergence in the accelerated
Collatz map $R_3(n) = (3n+1)/2^{\nu_2(3n+1)}$, and prove it as a necessary condition
\emph{conditional on} an explicit Local Equidistribution Hypothesis (LEH).
Defining $\fP$ as the empirical frequency of visits to the class $P = \{n \equiv 3 \pmod{4}\}$,
we show that, under LEH, any orbit diverging sub-exponentially must satisfy
$\fP \geq \fPstar = \logfour(8/3) = (3 - \logtwo 3)/2 \approx 0.7075$.
The unconditional part of the argument rests on the exact identity $\muP + \muN = -2$
between the ensemble-conditional expected drifts, under Haar measure, in the natural
metric $\Vmetric(n) = \log_{4/3}(n)$; LEH is the additional, unproven ingredient needed
to promote this ensemble-level identity to a per-orbit statement, and we relate it
explicitly to Chang's Weak Mixing Hypothesis \cite{chang2026paradigm}.
As corollaries we obtain:
\textup{(i)}~the exact $2$-adic drift $\PhiDrift(a) = \logtwo a - 2$,
recovering the unique contractive odd map $a=3$
(arithmetic isolation via $2^\phi \approx 3.0696$);
\textup{(ii)}~exact drift $-1$ in $\Vmetric$;
\textup{(iii)}~the Fibonacci structure of the first non-trivial convergents of $\logtwo 3$.
Empirical analysis of $25{,}000$ orbits ($n \leq 50{,}000$, odd) shows
maximum observed $\fP = 0.\overline{6}$, a $5.8\%$ margin below the threshold,
consistent with---but not proving---convergence.
We identify the gap between Tao's ``almost every'' \cite{tao2019} and ``every''
as the central open problem and propose two precise conjectures
(Universal Map Balance, Baire emptiness of the divergent subspace).
\end{abstract}

\maketitle

% ============================================
% 1. INTRODUCTION
% ============================================
\section{Introduction}\label{sec:intro}

The Collatz conjecture asks whether every positive integer $n$ eventually reaches
the cycle $(1,4,2)$ under iteration of
\begin{equation}\label{eq:collatz}
T(n) = \begin{cases}
n/2 & \text{if } n \text{ is even},\\
3n+1 & \text{if } n \text{ is odd}.
\end{cases}
\end{equation}
Despite verification for all $n < 2^{68}$ \cite{lagarias2010} and
profound structural results---most notably Tao's theorem that almost all orbits
attain almost bounded values \cite{tao2019}---the gap between ``almost every''
and ``every'' remains open.

This paper asks a sharper question: \emph{what exact condition would an orbit need
to satisfy in order to diverge?} We answer with a closed-form necessary condition
for sub-exponential divergence, proved conditional on an explicit hypothesis about
the orbit's local statistics (LEH, Hypothesis~\ref{hyp:leh}).
The condition is expressed in terms of the frequency $\fP$ with which an orbit
visits the residue class $P = \{n \equiv 3 \pmod{4}\}$, and the threshold is
\begin{equation}\label{eq:threshold}
\fPstar = \logfour\!\left(\frac{8}{3}\right) = \frac{3 - \logtwo 3}{2} \approx 0.7075.
\end{equation}
Under LEH, any orbit with $\fP < \fPstar$ is guaranteed to be contractive in expectation
(Theorem~\ref{thm:threshold}).

Our approach combines exact arithmetic results with empirical observation.
The threshold value itself is derived without any empirical input, from an identity
(Theorem~\ref{thm:identity}) that unifies the global drift $-1$ in the natural metric
$\Vmetric$ with the conditional drifts of the $P$- and $N$-class steps. Promoting that
ensemble-level identity into a necessary condition on a specific, deterministic orbit
is where LEH enters; without it, $\muP$ and $\muN$ are population averages over each
residue class, not guaranteed to equal the actual per-step drift experienced by any
one orbit's visited steps. The simulations in Section~\ref{sec:empirical} serve only
as corroboration and do not bear on LEH itself.

\subsection*{Main contributions}
\begin{enumerate}
\item \textbf{Theorem~\ref{thm:threshold}:} An exact closed-form necessary condition
for divergence, proved conditional on an explicit Local Equidistribution Hypothesis
(LEH, Hypothesis~\ref{hyp:leh}), which we relate to Chang's Weak Mixing Hypothesis.
\item \textbf{Theorem~\ref{thm:identity}:} The exact identity $\muP + \muN = -2$,
which algebraically unifies the global drift and the conditional drifts.
\item \textbf{Theorem~\ref{thm:drift2adic}:} The exact $2$-adic drift formula
$\PhiDrift(a) = \logtwo a - 2$, isolating $a=3$ as the unique contractive odd map.
\item \textbf{Theorem~\ref{thm:driftV}:} Exact drift $-1$ in the natural metric $\Vmetric$.
\item \textbf{Theorem~\ref{thm:fibonacci}:} The Fibonacci structure of the first
non-trivial convergents of $\logtwo 3$.
\item Strong empirical evidence from $25{,}000$ orbits, with a maximum observed $\fP$
remaining $5.8\%$ below the critical threshold.
\item Two open conjectures with precise gap analysis.
\end{enumerate}

\subsection*{Outline}
Section~\ref{sec:framework} sets up the mathematical framework.
Section~\ref{sec:main} presents the main result (Theorem~\ref{thm:threshold})
and its unifying identity (Theorem~\ref{thm:identity}).
Section~\ref{sec:supporting} collects the supporting rigorous results.
Section~\ref{sec:empirical} reports the simulation data.
Section~\ref{sec:conjectures} formulates the open conjectures.
Section~\ref{sec:gap} discusses the ``almost every'' vs ``every'' gap.
Section~\ref{sec:conclusions} concludes.

% ============================================
% 2. MATHEMATICAL FRAMEWORK
% ============================================
\section{Mathematical Framework}\label{sec:framework}

\subsection{The accelerated map}

Iterating $T$ directly produces long strings of even steps. It is standard to
accelerate the dynamics by absorbing all trailing divisions by $2$ into a single
macro step.

\begin{definition}[Accelerated Collatz map]\label{def:accelerated}
For odd $a \in \N$ and odd $n \in \N^+$,
\begin{equation}\label{eq:accelerated}
R_a(n) = \frac{an + 1}{2^{\nuTwo(an+1)}},
\end{equation}
where $\nuTwo(m)$ is the $2$-adic valuation of $m$ (the largest $k$ such that $2^k \mid m$).
\end{definition}

For $a=3$, $R_3$ is the standard accelerated Collatz map.
The orbit of $n$ under $R_3$ is the sequence of odd integers
$\Orb(n) = (n_0, n_1, n_2, \dots)$ with $n_0 = n$ and $n_{k+1} = R_3(n_k)$.
Convergence of the original map $T$ is equivalent to the existence of $K$ such
that $n_K = 1$.

\subsection{The $P/N$ classification}

The value of $\nuTwo(3n+1)$ depends only on $n \bmod 4$.

\begin{definition}[$P/N$ classes]\label{def:PN}
For odd $n$,
\begin{itemize}
\item $P = \{n \equiv 3 \pmod{4}\}$: here $\nuTwo(3n+1) = 1$ exactly, so one halving step occurs.
\item $N = \{n \equiv 1 \pmod{4}\}$: here $\nuTwo(3n+1) \geq 2$, so two or more halving steps occur.
\end{itemize}
\end{definition}

For an orbit of $K$ odd steps, the \emph{$P$-frequency} is
\begin{equation}\label{eq:fP}
\fP = \frac{|\{k : n_k \in P\}|}{K}.
\end{equation}

\subsection{Information metrics}

\begin{definition}[Logarithmic information]\label{def:info}
$I(n) = \logtwo n$.
\end{definition}

\begin{definition}[Natural metric of the inverse tree]\label{def:Vmetric}
$\Vmetric(n) = \log_{4/3} n = \dfrac{\ln n}{\ln(4/3)}$.
\end{definition}

The metric $\Vmetric$ is natural because the inverse tree of the Collatz map
has branching factor $\beta = 4/3$: every odd integer has exactly one even
preimage, and an odd preimage exists with asymptotic density $1/3$.

% ============================================
% 3. MAIN RESULT
% ============================================
\section{The Divergence Threshold}\label{sec:main}

\subsection{Conditional drifts}

Let $\muP$ and $\muN$ denote the expected change in $\Vmetric$ conditioned on
being in class $P$ and $N$, respectively.

\begin{lemma}[Conditional drifts]\label{lem:conditional}
For the Collatz map $R_3$,
\begin{equation}\label{eq:muP}
\muP = \frac{\logtwo 3 - 1}{2 - \logtwo 3} \approx +1.4094,
\end{equation}
\begin{equation}\label{eq:muN}
\muN = \frac{\logtwo 3 - 3}{2 - \logtwo 3} \approx -3.4094.
\end{equation}
\end{lemma}

\begin{proof}
For $n \in P$, we have $\nuTwo(3n+1) = 1$, so $R_3(n) = (3n+1)/2$.
Hence
\[
\muP = \log_{4/3}\!\left(\frac{3n+1}{2n}\right)
     = \frac{\logtwo(3/2)}{\logtwo(4/3)}
     = \frac{\logtwo 3 - 1}{2 - \logtwo 3}.
\]
For $n \in N$, we use the law of total expectation.
By Theorem~\ref{thm:drift2adic} (proved below), $\E[\nuTwo(3n+1)] = 2$ under
Haar measure on odd residues. Since $\PP(P) = \PP(N) = 1/2$,
\[
2 = \E[\nuTwo] = \tfrac{1}{2} \cdot 1 + \tfrac{1}{2} \cdot \E[\nuTwo \mid N],
\]
so $\E[\nuTwo \mid N] = 3$. Therefore
\[
\muN = \frac{\logtwo 3 - 3}{2 - \logtwo 3}. \qedhere
\]
\end{proof}

\subsection{The unifying identity}

\begin{theorem}[Algebraic unification]\label{thm:identity}
The conditional drifts satisfy the exact identity
\begin{equation}\label{eq:identity}
\muP + \muN = -2.
\end{equation}
\end{theorem}

\begin{proof}
Direct computation from Lemma~\ref{lem:conditional}:
\[
\muP + \muN
= \frac{(\logtwo 3 - 1) + (\logtwo 3 - 3)}{2 - \logtwo 3}
= \frac{2(\logtwo 3 - 2)}{2 - \logtwo 3}
= -2. \qedhere
\]
\end{proof}

\begin{remark}
With $f_P = 1/2$ (the Haar-measure expectation), the global drift is
$\tfrac{1}{2}\muP + \tfrac{1}{2}\muN = -1$,
recovering Theorem~\ref{thm:driftV} as a corollary of Theorem~\ref{thm:identity}.
\end{remark}

\subsection{A local equidistribution hypothesis}

The conditional drifts $\muP, \muN$ of Lemma~\ref{lem:conditional} are computed
as ensemble (Haar-measure) averages over the full population of odd integers in
each class. For $n \in P$, the per-step drift $\log_{4/3}((3n+1)/2n) \to \muP$ as
$n \to \infty$, so this substitution is asymptotically justified step-by-step.
For $n \in N$, however, $\nuTwo(3n+1)$ takes different finite values ($2,3,4,\dots$)
depending on $n$, and $\muN$ is only the \emph{average} of the resulting drift over
the whole class under Haar measure. Whether the $N$-steps actually visited by a
\emph{specific, deterministic} orbit reproduce that same average is a separate,
unproven question, which we isolate explicitly.

\begin{hypothesis}[Local Equidistribution Hypothesis, LEH]\label{hyp:leh}
For a given Collatz orbit $(n_0, n_1, n_2, \dots)$, let $S_N(K)$ denote the set of
indices $k < K$ with $n_k \in N$. Then
\[
\frac{1}{|S_N(K)|} \sum_{k \in S_N(K)} \nuTwo(3n_k+1) \;\longrightarrow\; 3
\qquad \text{as } K \to \infty.
\]
\end{hypothesis}

\begin{remark}
LEH is not proven for any individual orbit, and we do not prove it here. Proving it
for every orbit is, in conjunction with Theorem~\ref{thm:threshold}, at least as hard
as the Collatz conjecture itself (see Conjecture~\ref{conj:balance} and the discussion
in Section~\ref{sec:gap}). LEH is the same kind of statement as Chang's Weak Mixing
Hypothesis \cite{chang2026paradigm}: both require a deterministic orbit's local time
average to match an ensemble (Haar) average of a genuinely orbit-dependent quantity.
We state LEH explicitly so that Theorem~\ref{thm:threshold} below is exactly as strong
as its proof, and no stronger.
\end{remark}

\subsection{The divergence threshold}

\begin{theorem}[Divergence threshold, conditional on LEH]\label{thm:threshold}
Assume Hypothesis~\ref{hyp:leh} (LEH) holds for a given Collatz orbit. Then a
necessary condition for sub-exponential divergence of that orbit is
\begin{equation}\label{eq:threshold-thm}
\fP \geq \fPstar = \frac{3 - \logtwo 3}{2} = \logfour\!\left(\frac{8}{3}\right) \approx 0.7075187.
\end{equation}
Equivalently, under LEH, any orbit with $\fP < \fPstar$ is guaranteed to be
contractive in expectation.
\end{theorem}

\begin{proof}
For sub-exponential divergence, the time-averaged drift in $\Vmetric$ must satisfy
$\frac{1}{K}\sum_{k=0}^{K-1} \Delta \Vmetric(n_k) \to 0$.
Split the sum over visited $P$-steps and $N$-steps:
\[
\frac{1}{K}\sum_{k=0}^{K-1} \Delta \Vmetric(n_k)
= f_P \cdot \Bigl[\text{avg.\ drift on visited } P\text{-steps}\Bigr]
+ (1-f_P) \cdot \Bigl[\text{avg.\ drift on visited } N\text{-steps}\Bigr].
\]
The bracketed $P$-term converges to $\muP$ unconditionally, since the per-step
drift $\to \muP$ pointwise as $n \to \infty$ (Lemma~\ref{lem:conditional}). The
bracketed $N$-term converges to $\muN$ \emph{precisely because} LEH guarantees
that the visited $N$-steps' average $\nuTwo$ converges to $3$, the same value
that defines $\muN$ under Haar measure; without LEH this term need not equal
$\muN$ for a specific orbit. Under LEH, then, the time-averaged drift converges to
the convex combination
\[
f_P \cdot \muP + (1 - f_P) \cdot \muN,
\]
and non-negative drift requires
\[
f_P \cdot \muP + (1 - f_P) \cdot \muN \geq 0,
\]
i.e.
\[
f_P \geq \frac{-\muN}{\muP - \muN}.
\]
Substituting the exact values from Lemma~\ref{lem:conditional}:
\[
\frac{-\muN}{\muP - \muN}
= \frac{(3 - \logtwo 3)/(2 - \logtwo 3)}{2/(2 - \logtwo 3)}
= \frac{3 - \logtwo 3}{2}
= \logfour\!\left(\frac{8}{3}\right). \qedhere
\]
\end{proof}

\begin{corollary}
If every Collatz orbit satisfies LEH and $\fP \to 1/2$ as $K \to \infty$,
then no orbit diverges (since $1/2 < \fPstar$). Conjecture~\ref{conj:balance}
(Universal Map Balance) and LEH are two faces of the same underlying orbit-level
equidistribution phenomenon; we do not claim one implies the other, only that a
proof of Collatz along this route would need both.
\end{corollary}

% ============================================
% 4. SUPPORTING RESULTS
% ============================================
\section{Supporting Results}\label{sec:supporting}

\subsection{Exact $2$-adic drift}

\begin{theorem}[$2$-adic drift formula]\label{thm:drift2adic}
For the accelerated map $R_a$ with odd $a$, under Haar measure uniform over
odd residues modulo $2^K$,
\begin{equation}\label{eq:drift2adic}
\PhiDrift(a) = \E[\logtwo(R_a(n)/n)] = \logtwo a - 2.
\end{equation}
\end{theorem}

\begin{proof}
Because $\gcd(a,2)=1$, the integer $a$ is invertible modulo $2^k$.
The condition $\nuTwo(an+1) \geq k$ is equivalent to
$n \equiv -a^{-1} \pmod{2^k}$.
Since $a^{-1}$ is odd, $-a^{-1} \equiv 1 \pmod{2}$, so the required residue
is always odd and well-defined.
Among the $2^{k-1}$ odd residue classes modulo $2^k$, exactly one satisfies
this condition. Hence
\[
\PP(\nuTwo(an+1) \geq k \mid n \text{ odd}) = \frac{1}{2^{k-1}} \quad (k \geq 1).
\]
Therefore
\[
\E[\nuTwo(an+1)] = \sum_{k=1}^{\infty} \PP(\nuTwo(an+1) \geq k)
                 = \sum_{k=1}^{\infty} \frac{1}{2^{k-1}}
                 = 2.
\]
For large $n$, $\logtwo((an+1)/n) \to \logtwo a$, so
\[
\PhiDrift(a) = \logtwo a - \E[\nuTwo(an+1)] = \logtwo a - 2. \qedhere
\]
\end{proof}

\begin{corollary}\label{cor:drift-values}
\begin{center}
\begin{tabular}{@{}ccc@{}}
\toprule
$a$ & $\PhiDrift(a)$ & Behavior \\
\midrule
$1$ & $-2.000$ & trivial collapse \\
$3$ & $-0.415$ & contractive \\
$5$ & $+0.322$ & expansive \\
$7$ & $+0.807$ & explosive \\
\bottomrule
\end{tabular}
\end{center}
\end{corollary}

\subsection{Exact drift in the natural metric}

\begin{theorem}[Drift in $\Vmetric$]\label{thm:driftV}
For $a=3$,
\begin{equation}\label{eq:driftV}
\E[\Delta \Vmetric] = -1 \quad \text{per macro step}.
\end{equation}
\end{theorem}

\begin{proof}
From Theorem~\ref{thm:drift2adic} with $a=3$,
\[
\E[\Delta \Vmetric]
= \frac{\E[\logtwo(R_3(n)/n)]}{\logtwo(4/3)}
= \frac{\logtwo 3 - 2}{2 - \logtwo 3}
= -1. \qedhere
\]
\end{proof}

\subsection{Arithmetic isolation of $a=3$}

\begin{theorem}[Isolation]\label{thm:isolation}
Among all odd integers $a > 1$:
\begin{enumerate}
\item[(i)] $R_a$ is contractive in expectation ($\PhiDrift(a) < 0$) if and only if $a < 4$,
i.e.~$a = 3$.
\item[(ii)] The unique odd integer in $(1, 2^\phi)$, where $\phi = (1+\sqrt{5})/2$,
is $a = 3$ (since $2^\phi \approx 3.0696$).
\end{enumerate}
\end{theorem}

\begin{proof}
Part~(i) follows from Theorem~\ref{thm:drift2adic}:
$\PhiDrift(a) < 0 \iff \logtwo a < 2 \iff a < 4$.
Among odd integers $>1$, only $a=3$ satisfies this.

Part~(ii) is arithmetic: $2^\phi \approx 3.0696$, and $3 < 2^\phi < 5$.
\end{proof}

\begin{remark}
Part~(i) is dynamically essential: the transition from contractive to expansive
occurs at $a=4$, and $a=3$ is the only non-trivial odd integer on the
contractive side. Part~(ii) is an independent arithmetic observation; we do not
claim that the golden ratio is causally responsible for Collatz convergence.
\end{remark}

\subsection{Fibonacci structure of $\logtwo 3$ convergents}

\begin{theorem}[Fibonacci convergents]\label{thm:fibonacci}
The continued fraction of $\logtwo 3$ begins
\[
\logtwo 3 = [1; 1, 1, 2, 2, 3, 1, 5, 2, 23, \dots].
\]
The first two non-trivial convergents are ratios of consecutive Fibonacci numbers:
\[
\frac{3}{2} = \frac{F_4}{F_3}, \qquad \frac{8}{5} = \frac{F_6}{F_5}.
\]
\end{theorem}

\begin{proof}
Direct computation from the recursive convergent formula.
With $a_0=1, a_1=1, a_2=1$ we obtain $h_2/k_2 = 3/2$;
with $a_3=2$ we obtain $h_3/k_3 = 8/5$.
\end{proof}

\begin{remark}
This structure is a consequence of the initial CF coefficients $[1;1,1,\dots]$
matching those of $\phi$. The pattern breaks at the fourth coefficient.
The theorem is an arithmetic observation; its dynamical significance beyond
motivating Baker-type cycle bounds is an open question.
\end{remark}

% ============================================
% 5. EMPIRICAL EVIDENCE
% ============================================
\section{Empirical Evidence}\label{sec:empirical}

We simulated all orbits for odd starting points $n \leq 50{,}000$ under the
accelerated map $R_3$, with a per-orbit cap of $100{,}000$ steps. Every orbit
reached $1$ well within the cap (none were truncated), so every $K$ and $\fP$
reported below is a true, full-orbit value rather than a value computed on a
truncated prefix. For each orbit, $K$ is the number of odd steps until
reaching $1$, and $\fP$ is the fraction of those steps belonging to class $P$.

\begin{table}[ht]
\centering
\caption{Simulation summary ($25{,}000$ orbits, odd $n \leq 50{,}000$)}
\label{tab:simulation}
\begin{tabular}{@{}lr@{}}
\toprule
Metric & Value \\
\midrule
Orbits analyzed & $25{,}000$ \\
Orbits truncated by step limit & $0$ \\
Mean $\fP$ per orbit & $0.4561$ \\
Max $\fP$ observed & $0.\overline{6}$ ($n = 151$ and $n = 1431$) \\
Orbits with $\fP \geq \fPstar$ & $0$ ($0\%$) \\
Empirical margin below threshold ($\fPstar - \max \fP$) & $0.0409$ ($5.8\%$) \\
Lag-$1$ autocorr.~of $P/N$ steps & $-0.0045$ \\
\bottomrule
\end{tabular}
\end{table}

\textbf{No orbit approaches the threshold.} Within this finite, empirically
checked range, the maximum observed $\fP = 2/3$ remains $5.8\%$ below the
critical value $\fPstar \approx 0.7075$. This is a statement about the sample
analyzed, not a general bound on how close any orbit could come to $\fPstar$.

The per-orbit mean $\fP = 0.456$ differs from the step-weighted mean
$\approx 0.500$ because short orbits---which tend to start close to $1$ and
have fewer $P$-class steps---contribute equally in the per-orbit average.
Both measures are consistent with $\fP \to 1/2$ for long orbits.

The near-zero lag-$1$ autocorrelation ($\bar{\rho}_1 \approx -0.005$) of the
binary $P/N$ sequence suggests approximate independence of successive steps.
This is consistent with a random-walk model in which $\fP$ concentrates
around $1/2$ by the Law of Large Numbers, with variance decaying as $1/K$.

% ============================================
% 6. OPEN CONJECTURES
% ============================================
\section{Open Conjectures}\label{sec:conjectures}

\subsection{Universal Map Balance}

\begin{conjecture}[Universal Map Balance]\label{conj:balance}
For every Collatz orbit, the empirical frequency satisfies $\fP \to 1/2$
as the orbit length grows.
\end{conjecture}

\textbf{What this would imply.} By Theorem~\ref{thm:threshold}, if LEH
(Hypothesis~\ref{hyp:leh}) holds and $\fP < \fPstar$ for every orbit, then no
orbit diverges. Conjecture~\ref{conj:balance} together with LEH is therefore
\emph{sufficient} (via Theorem~\ref{thm:threshold}) to prove the Collatz
conjecture; neither piece is proved here.

\textbf{Known results.} Tao \cite{tao2019} proved that almost all orbits
(logarithmic density $1$) attain almost bounded values---a distributional result.
Chang \cite{chang2026reduction} proved the Map Balance Theorem at the map level:
among burst residues modulo $2^K$, the counts mapping to gap-start classes
$\equiv 3$ vs $\equiv 7 \pmod{8}$ differ by exactly $1$ for all $K \geq 5$.
This shows that \emph{all residual imbalance is orbit-level, not map-level},
reducing the conjecture to whether every individual orbit visits two residue
classes modulo $32$ with sufficient balance.

\subsection{Baire emptiness of the divergent subspace}

\begin{conjecture}[$\Sigma_{\mathrm{div}} = \emptyset$]\label{conj:baire}
The subspace of divergent orbits in the shift space $\Sigma_{\mathrm{Collatz}}$ is empty.
\end{conjecture}

A potential argument would require:
\begin{enumerate}
\item[(i)] $\Sigma_{\mathrm{Collatz}}$ is a complete metric (Baire) space;
\item[(ii)] $\Sigma_{\mathrm{div}}$ is of first category;
\item[(iii)] the shift $\sigma_{\mathrm{Collatz}}$ is topologically mixing.
\end{enumerate}
None of (i)--(iii) has been verified. This conjecture is presented as a
potential topological pathway, not as a near-complete argument.

% ============================================
% 7. THE GAP
% ============================================
\section{The Gap: ``Almost Every'' vs ``Every''}\label{sec:gap}

\subsection{What is known}

Table~\ref{tab:known} summarizes the rigorous and empirical results established
in this paper and by prior authors.

\begin{table}[ht]
\centering
\caption{Established results}
\label{tab:known}
\begin{tabular}{@{}ll@{}}
\toprule
Result & Status \\
\midrule
$\PhiDrift(a) = \logtwo a - 2$ (Theorem~\ref{thm:drift2adic}) & Proven \\
$\E[\Delta \Vmetric] = -1$ (Theorem~\ref{thm:driftV}) & Proven \\
$a=3$ unique contractive odd $a$ among $R_a$, $a>1$ (Theorem~\ref{thm:isolation}) & Proven \\
CF convergents of $\logtwo 3$ are Fibonacci ratios (Theorem~\ref{thm:fibonacci}) & Proven \\
$\fP \geq \logfour(8/3)$ necessary for divergence (Theorem~\ref{thm:threshold}) & Proven, conditional on LEH \\
$\muP + \muN = -2$ (Theorem~\ref{thm:identity}) & Proven \\
$25{,}000$ orbits: none exceed threshold & Empirical \\
Autocorrelation of $P/N$ steps $\approx 0$ & Empirical \\
Map-level balance: residual bias is orbit-level & Chang \cite{chang2026reduction} \\
Almost all orbits attain almost bounded values & Tao \cite{tao2019} \\
\bottomrule
\end{tabular}
\end{table}

\subsection{What remains open}

\begin{table}[ht]
\centering
\caption{Open problems}
\label{tab:open}
\begin{tabular}{@{}ll@{}}
\toprule
Question & Difficulty \\
\midrule
LEH for every orbit (Hypothesis~\ref{hyp:leh}) & Very high \\
$\fP \to 1/2$ for \emph{every} orbit (Conjecture~\ref{conj:balance}) & Very high \\
$\Sigma_{\mathrm{div}} = \emptyset$ via Baire (Conjecture~\ref{conj:baire}) & Extreme \\
The Collatz conjecture & Unknown \\
\bottomrule
\end{tabular}
\end{table}

\subsection{The ergodic pathway}

The empirical observation that $P/N$ steps have near-zero autocorrelation
suggests an ergodic approach to Conjecture~\ref{conj:balance}.
If the steps were exactly independent with $\PP(P) = 1/2$,
then by Hoeffding's inequality
\[
\PP\bigl(\fP \geq \fPstar \mid \text{orbit of length } K\bigr)
\leq \exp\!\bigl(-2K(\fPstar - 1/2)^2\bigr).
\]
For $K = 1000$, this bound is approximately $10^{-38}$.
We emphasize that this calculation \emph{assumes independence},
which is precisely what would need to be proved.
It is offered as motivation, not as an argument that divergence is unlikely.

Chang's companion paper \cite{chang2026paradigm} establishes a
\emph{Paradigm Exhaustion Theorem} across $29$ distinct mathematical frameworks:
every known approach to promoting distributional convergence to pointwise
convergence encounters an irreducible structural obstruction at the orbit level.
This provides independent, systematic evidence that the gap between
``almost every'' and ``every'' is not a technical inconvenience but a
fundamental structural barrier.

The critical open question is: \emph{can the approximate independence observed
empirically be elevated to a rigorous mixing statement?}

% ============================================
% 8. CONCLUSIONS
% ============================================
\section{Conclusions}\label{sec:conclusions}

We have \emph{not} proved the Collatz conjecture.
What we have proved is a rigorous structural characterization of why the
conjecture is arithmetically special and where the exact boundary of divergence lies:

\begin{enumerate}
\item The exact contractive drift $\PhiDrift(3) = \logtwo 3 - 2 < 0$,
unique among odd $a > 1$ (Theorem~\ref{thm:drift2adic}).
\item The exact drift $-1$ in the natural metric $\Vmetric$ (Theorem~\ref{thm:driftV}).
\item The exact closed-form divergence threshold
$\fPstar = \logfour(8/3)$, proved conditional on the Local Equidistribution
Hypothesis (Theorem~\ref{thm:threshold}, Hypothesis~\ref{hyp:leh}).
\item The unifying identity $\muP + \muN = -2$ (Theorem~\ref{thm:identity}),
from which Theorem~\ref{thm:driftV} follows as a corollary.
\item Strong empirical evidence: $25{,}000$ orbits, none approaching the threshold.
\end{enumerate}

The gap between ``almost every'' (Tao \cite{tao2019}) and ``every'' persists.
Chang \cite{chang2026paradigm} demonstrates that this gap is not a failure of
any particular method but a structural feature of the problem.
The conjecture reduces to whether every individual orbit achieves the balance
predicted by the map-level structure---a one-bit, pointwise mixing question.

Closing the gap requires proving that $\fP \to 1/2$ for every individual orbit,
a statement that ergodic theory reaches only on average.
The near-zero autocorrelation of $P/N$ steps is an empirical hint that a mixing
argument may be within reach, but the proof remains open.

% ============================================
% REFERENCES
% ============================================
\begin{thebibliography}{99}

\bibitem{tao2019}
T.~Tao, \textit{Almost all orbits of the Collatz map attain almost bounded values},
arXiv:1909.03562 [math.PR], 2019.

\bibitem{chang2026reduction}
E.~Y.~Chang, \textit{A structural reduction of the Collatz conjecture to one-bit orbit mixing},
arXiv:2603.25753 [math.DS], Stanford University, 2026.

\bibitem{chang2026paradigm}
E.~Y.~Chang, \textit{Exploring Collatz dynamics with human--LLM collaboration},
arXiv:2603.11066 [math.DS], Stanford University, 2026.

\bibitem{lagarias2010}
J.~C.~Lagarias, \textit{The Ultimate Challenge: The $3x+1$ Problem},
American Mathematical Society, 2010.

\end{thebibliography}

\end{document}
